%% ============================================================
%%  Presentation Script — Pneumonia Classification
%%  Carolina Reis & Jakub Błaszczyk — University of Aveiro
%%  10 min presentation + 3 min Q&A
%%  Color coding:
%%    \jakub{} — dark blue/teal box  → Jakub speaks
%%    \carol{} — rose/pink box       → Carolina speaks
%% ============================================================
\documentclass[12pt,a4paper]{article}

\usepackage[utf8]{inputenc}
\usepackage[T1]{fontenc}
\usepackage[margin=2cm]{geometry}
\usepackage{xcolor}
\usepackage{tcolorbox}
\tcbuselibrary{skins,breakable}
\usepackage{titlesec}
\usepackage{enumitem}
\usepackage{parskip}
\usepackage{amsmath}
\usepackage{booktabs}
\usepackage{array}
\usepackage{hyperref}

%% ---- Color palette (matches slide theme) ----
\definecolor{RD}{RGB}{130,25,70}       % dark rose — slide accent
\definecolor{RM}{RGB}{195,65,110}      % mid rose
\definecolor{RL}{RGB}{253,234,244}     % light rose background
\definecolor{JB}{RGB}{20,60,120}       % Jakub — dark blue
\definecolor{JBL}{RGB}{210,225,250}    % Jakub — light blue
\definecolor{CR}{RGB}{130,25,70}       % Carolina — dark rose (= RD)
\definecolor{CRL}{RGB}{253,234,244}    % Carolina — light rose (= RL)
\definecolor{GLOSS}{RGB}{45,100,80}    % Glossary — dark teal
\definecolor{GLOSSL}{RGB}{210,240,230} % Glossary — light teal
\definecolor{QA}{RGB}{90,40,120}       % Q&A — purple
\definecolor{QAL}{RGB}{235,220,250}    % Q&A — light purple
\definecolor{SUM}{RGB}{40,40,40}       % Summary — near-black

%% ---- Speaker box styles ----
\tcbset{
  jakubbox/.style={
    breakable,
    colback=JBL, colframe=JB,
    fonttitle=\bfseries\small\color{white},
    colbacktitle=JB,
    boxrule=1pt, arc=5pt,
    left=8pt, right=8pt, top=5pt, bottom=5pt,
    toptitle=3pt, bottomtitle=3pt,
    title={#1}
  },
  carolbox/.style={
    breakable,
    colback=CRL, colframe=CR,
    fonttitle=\bfseries\small\color{white},
    colbacktitle=CR,
    boxrule=1pt, arc=5pt,
    left=8pt, right=8pt, top=5pt, bottom=5pt,
    toptitle=3pt, bottomtitle=3pt,
    title={#1}
  },
  glossbox/.style={
    breakable,
    colback=GLOSSL, colframe=GLOSS,
    fonttitle=\bfseries\small\color{white},
    colbacktitle=GLOSS,
    boxrule=1pt, arc=5pt,
    left=8pt, right=8pt, top=5pt, bottom=5pt,
    toptitle=3pt, bottomtitle=3pt,
    title={#1}
  },
  qabox/.style={
    breakable,
    colback=QAL, colframe=QA,
    fonttitle=\bfseries\small\color{white},
    colbacktitle=QA,
    boxrule=1pt, arc=5pt,
    left=8pt, right=8pt, top=5pt, bottom=5pt,
    toptitle=3pt, bottomtitle=3pt,
    title={#1}
  },
  sumbox/.style={
    breakable,
    colback=gray!8, colframe=SUM,
    fonttitle=\bfseries\small\color{white},
    colbacktitle=SUM,
    boxrule=1pt, arc=5pt,
    left=8pt, right=8pt, top=5pt, bottom=5pt,
    toptitle=3pt, bottomtitle=3pt,
    title={#1}
  },
}

\newcommand{\jakub}[2]{%
  \begin{tcolorbox}[jakubbox={JAKUB \textbullet\ #1}]
  #2
  \end{tcolorbox}
}
\newcommand{\carol}[2]{%
  \begin{tcolorbox}[carolbox={CAROLINA \textbullet\ #1}]
  #2
  \end{tcolorbox}
}

%% ---- Section headings ----
\titleformat{\section}[block]
  {\large\bfseries\color{RD}}
  {\thesection.}{0.5em}{}[\color{RD}\rule{\linewidth}{0.8pt}]

\titleformat{\subsection}[block]
  {\normalsize\bfseries\color{JB}}
  {\thesubsection.}{0.5em}{}

%% ---- Slide label command ----
\newcommand{\slidetag}[1]{%
  \noindent{\small\color{gray}\textit{[Slide #1]}}\par\smallskip
}

%% ---- Timing note command ----
\newcommand{\timing}[1]{%
  \noindent{\footnotesize\color{gray!70!black}\textbf{Timing:} #1}\par\smallskip
}

%% ============================================================
\begin{document}
%% ============================================================

\begin{center}
  {\LARGE\bfseries\color{RD} Presentation Script}\\[0.4em]
  {\large Pneumonia Classification from Pediatric Chest X-Rays}\\[0.2em]
  {\normalsize Carolina Reis \& Jakub B\l{}aszczyk \quad$\cdot$\quad University of Aveiro, DETI \quad$\cdot$\quad June 2026}\\[0.3em]
  {\small\color{gray} Total time: \textbf{10 min} presentation $+$ \textbf{3 min} Q\&A\quad|\quad
  \textcolor{JB}{\rule{1.2em}{0.8em}}~Jakub \quad
  \textcolor{CR}{\rule{1.2em}{0.8em}}~Carolina}
\end{center}

\vspace{0.5em}
\noindent\rule{\linewidth}{1.2pt}
\vspace{0.3em}

%% ============================================================
\section*{INTRO — Jakub (approx.\ 2 min 30 s)}
%% ============================================================

% ------ SLIDE 1 ------
\slidetag{1 — Cover}
\timing{$\sim$20 s}

\jakub{Slide 1 — Cover}{
Good morning, everyone. Today we're presenting our work on pneumonia classification from pediatric chest X-rays using lightweight deep learning models. The core question we set out to answer is: can a small, efficient neural network reliably detect pneumonia from X-rays --- and can we trust what it's looking at?
}

\vspace{0.4em}

% ------ SLIDE 2 ------
\slidetag{2 — Table of Contents (Part 1)}
\timing{$\sim$20 s}

\jakub{Slide 2 — Table of Contents}{
Let me quickly walk you through the structure of our presentation. We'll start with the Outline, then move into Motivation and the clinical problem, the Dataset, Preprocessing, and the Pipeline and Model Architectures. From there we'll cover the Training Setup and then hand over to Carolina for the results --- Quantitative Results, a Visual Comparison of the models, and the Training Curves for our best model, DenseNet121.
}

\vspace{0.4em}

% ------ SLIDE 3 ------
\slidetag{3 — Table of Contents (Part 2)}
\timing{$\sim$20 s}

\jakub{Slide 3 — Table of Contents continued}{
Carolina continues with the Confusion Matrix for the binary ensemble, our two innovations --- Three-Class Classification and the Hierarchical Two-Stage Classifier --- and the Ensemble and Calibration Analysis. I then take over for the three Grad-CAM explainability slides. And Carolina closes with Ethics and Limitations, and the Conclusions.
}

\vspace{0.4em}

% ------ SLIDE 4 ------
\slidetag{4 — Outline}
\timing{$\sim$25 s}

\jakub{Slide 4 — Outline}{
Structurally, the work flows through seven sections. We begin with motivation and the problem definition, move into dataset and preprocessing, then models and training. We then present our binary classification results, followed by our main innovation --- a three-class extension distinguishing Normal, Bacterial, and Viral pneumonia, as well as a hierarchical two-stage variant of that classifier. We then apply Grad-CAM for explainability, and close with ethics and conclusions. The central highlighted box --- Innovations --- is where we go beyond the standard binary task, so that's the part I'd invite you to pay particular attention to.
}

\vspace{0.4em}

% ------ SLIDE 5 ------
\slidetag{5 — Motivation \& Problem Statement}
\timing{$\sim$35 s}

\jakub{Slide 5 — Motivation \& Problem Statement}{
Pneumonia kills over 2.5 million people per year and is the leading cause of death in children under five. Diagnosis typically requires a trained radiologist reading a chest X-ray --- a resource that simply doesn't exist in many low-income settings. Late or missed diagnosis translates directly into preventable deaths. X-rays are the most accessible imaging modality globally, which is exactly why automating their interpretation could have enormous impact. And as you can see in the three X-rays on the left, the visual difference between a normal lung, a bacterial pneumonia, and a viral pneumonia can be very subtle --- which makes this a genuinely hard problem.

Our scope covers binary and three-class classification, Grad-CAM explainability, ensemble methods, and calibration. We explicitly exclude segmentation and clinical deployment --- this is a research prototype.
}

\vspace{0.4em}

% ------ SLIDE 6 ------
\slidetag{6 — Dataset: Kermany Chest X-Ray Images}
\timing{$\sim$35 s}

\jakub{Slide 6 — Dataset}{
We used the publicly available Kermany dataset, sourced from Guangzhou Women and Children's Medical Center and published on Kaggle in 2018. It contains 5,863 pediatric chest X-rays labelled as either NORMAL or PNEUMONIA. Crucially, the filenames encode a finer distinction --- BACTERIA or VIRUS --- which we leveraged to create three-class labels without any additional annotation cost. The dataset is heavily imbalanced: 74\% of training images are pneumonia cases. This imbalance directly influenced our choice of evaluation metrics --- we prioritised recall and AUROC over raw accuracy.
}

\vspace{0.4em}

% ------ SLIDE 7 ------
\slidetag{7 — Preprocessing \& Data Augmentation}
\timing{$\sim$35 s}

\jakub{Slide 7 — Preprocessing \& Data Augmentation}{
All images were resized to 224$\times$224 pixels to match the input standard for ResNet and DenseNet. Since the X-rays are grayscale, we replicated the single channel three times to produce an RGB image compatible with ImageNet-pretrained backbones. We then normalized using ImageNet statistics. For the training split only, we applied light augmentation --- random horizontal flips and rotations up to $\pm$10 degrees. We kept augmentation conservative deliberately: aggressive geometric transforms could distort clinically meaningful patterns like opacities and consolidations.
}

\vspace{0.4em}

% ------ SLIDE 8 ------
\slidetag{8 — Pipeline \& Model Architectures}
\timing{$\sim$40 s}

\jakub{Slide 8 — Pipeline \& Model Architectures}{
Our pipeline takes a chest X-ray, resizes and normalizes it, passes it through a feature extractor, applies a sigmoid or softmax output, and then evaluates with both metrics and Grad-CAM. We tested three architectures. The custom CNN is our baseline --- three convolutional blocks with no pretrained weights and roughly 0.5 million parameters. ResNet18 brings skip connections and 11 million ImageNet-pretrained parameters. DenseNet121 is our primary model of interest --- it has dense connectivity between all layers, just 7 million parameters, and carries the CheXNet heritage, having been used by Rajpurkar et al.\ in 2017 to match radiologist-level performance on chest X-rays. Its dense feature reuse makes it particularly good at detecting subtle diffuse opacities.
}

\vspace{0.4em}

% ------ SLIDE 9 ------
\slidetag{9 — Training Setup}
\timing{$\sim$35 s}

\jakub{Slide 9 — Training Setup}{
All models shared the same hyperparameter configuration: Adam optimizer with a learning rate of $10^{-4}$, ReduceOnPlateau scheduling, batch size of 32, early stopping with patience 5, and a maximum of 30 epochs. We ran each experiment three times with different random seeds to ensure robustness. For the pretrained models, we used a two-phase strategy: Phase 1 freezes the backbone and trains only the classification head for 5 epochs, giving the new head a chance to stabilize before Phase 2, where the entire network is unfrozen for full fine-tuning. Evaluation was done strictly on a fixed test set, evaluated once at the very end.
}

%% ============================================================
\section*{ACTUAL CONTENT — Carolina (approx.\ 5 min)}
%% ============================================================

% ------ SLIDE 10 ------
\slidetag{10 — Quantitative Results: Mean $\pm$ Std (3 seeds)}
\timing{$\sim$45 s}

\carol{Slide 10 — Quantitative Results}{
Here are the core results across three seeds. DenseNet121 achieves a recall of 0.997 --- meaning it missed only 1 to 2 pneumonia cases per run out of 390. Its AUROC of 0.966 actually exceeds the CheXNet-14 reference of 0.927. Recall was our primary success criterion, and all three models exceed 0.90 --- but the transfer learning models are clearly superior. The custom CNN shows high variance, especially in specificity, reflecting the instability of training from scratch. And importantly, inference time is essentially identical across all three models at around 13 to 15 milliseconds per image --- meaning there is no computational cost to choosing the better model.
}

\vspace{0.4em}

% ------ SLIDE 11 ------
\slidetag{11 — Visual Comparison of Models}
\timing{$\sim$35 s}

\carol{Slide 11 — Visual Comparison of Models}{
This chart confirms what the numbers tell us. F1 and accuracy are similar between DenseNet121 and ResNet18, both far ahead of the custom CNN. Recall is above 0.90 for all models, with transfer learning models approaching 1.0. The most striking gap is in AUROC, where transfer learning clearly dominates.

One important trade-off to note: maximising recall inevitably reduces specificity. This means some healthy patients will be flagged as potential pneumonia cases. In a triage context, this is acceptable --- a clinician reviews flagged results. It is not acceptable to miss true cases. The CNN specificity of $0.49 \pm 0.36$ is particularly unstable across seeds, reflecting the variance of training from scratch.
}

\vspace{0.4em}

% ------ SLIDE 12 ------
\slidetag{12 — Training Curves: DenseNet121 (Best Model)}
\timing{$\sim$35 s}

\carol{Slide 12 — Training Curves: DenseNet121}{
The training curves show two distinct phases. In Phase 1, the loss drops rapidly --- the ImageNet features are already strong, so the new head converges quickly. In Phase 2, we see gradual, stable improvement as the full network is fine-tuned. The learning rate scheduler reduces on plateau, and early stopping prevents overfitting. Across all three seeds, the standard deviation on all metrics is approximately 0.020 --- indicating a robust, low-variance training procedure.
}

\vspace{0.4em}

% ------ SLIDE 13 ------
\slidetag{13 — Confusion Matrix: Binary Ensemble (6 checkpoints)}
\timing{$\sim$35 s}

\carol{Slide 13 — Confusion Matrix: Binary Ensemble}{
This confusion matrix shows results for our binary ensemble, which combines predictions from six checkpoints across the three seeds of DenseNet121 and ResNet18. Out of 390 true pneumonia cases, only 1 was missed --- a recall of 0.997. There are 89 false positives, giving a specificity of 0.620 and precision of 0.814. Clinically, the 89 false alarms are not a failure --- they are flagged for clinician review. The single missed case is what matters, and the ensemble nearly eliminates that risk entirely.
}

\vspace{0.4em}

% ------ SLIDE 14 ------
\slidetag{14 — Innovation I: Three-Class Classification}
\timing{$\sim$45 s}

\carol{Slide 14 — Innovation I: Three-Class Classification}{
Beyond binary detection, we extended the task to distinguish Normal, Bacterial, and Viral pneumonia --- a clinically meaningful distinction since bacterial infections typically require antibiotics while viral ones do not. The labels were derived automatically from the filenames, at no additional annotation cost, using a patient-aware split to prevent leakage.

DenseNet121 achieves 0.799 accuracy and 0.784 macro-F1 on this harder task. The ensemble of ResNet18 and DenseNet121 across three seeds pushes this to 0.822 accuracy and 0.808 macro-F1, with an AUROC of 0.960. Looking at the per-class breakdown and the confusion matrix: Bacteria is well classified --- recall 0.97, F1 0.92. Virus, however, is the hardest class, with F1 only 0.72. The dominant error is Normal being misclassified as Virus, because viral infiltrates can be faint and diffuse, resembling a clear lung. This is a data-level bottleneck, not a modelling failure, and the three-class signal in the filenames is genuinely learnable with the same lightweight architecture.
}

\vspace{0.4em}

% ------ SLIDE 15 ------
\slidetag{15 — Innovation II: Hierarchical Two-Stage Classifier}
\timing{$\sim$40 s}

\carol{Slide 15 — Innovation II: Hierarchical Two-Stage Classifier}{
Our second innovation is a hierarchical two-stage classifier, motivated by the original Kermany paper --- which, remarkably, \textbf{never trained a flat three-way classifier}, but rather used two separate binary classifiers. We replicate that design philosophy and compare it head-to-head with our flat softmax approach.

The architecture is shown in the diagram. Stage A is the binary DenseNet121 classifier we already trained: it decides Normal versus Pneumonia. If the output is Pneumonia, the image is passed to Stage B --- a separate DenseNet121 trained only on pneumonia cases, which decides Bacteria versus Virus. Normal cases never reach Stage B.

The results table compares flat versus hierarchical. The flat model achieves F1 macro 0.784; the hierarchical model achieves 0.774 --- \textbf{the hierarchical approach does not win}. Why? Stage A errors propagate to Stage B with no possibility of recovery --- a Normal case misclassified as Pneumonia in Stage A will be further misclassified in Stage B. The Normal-to-Virus confusion is unchanged by the cascade. The flat softmax weighs all three classes jointly and is more robust.

The result is \textbf{negative but informative}: the bottleneck is intrinsic to the data, not to the modelling approach. Virus samples genuinely look like Normal to both architectures.
}

\vspace{0.4em}

% ------ SLIDE 16 ------
\slidetag{16 — Ensemble \& Calibration Analysis}
\timing{$\sim$40 s}

\carol{Slide 16 — Ensemble \& Calibration Analysis}{
Let me now cover two additional analyses: ensembling and calibration.

On ensembling: our binary ensemble averages probabilities from six checkpoints --- DenseNet121 and ResNet18 across three seeds. It achieves F1 0.896 and recall 0.997, missing only one case out of 390. For the three-class task, the ensemble is the best model on every macro metric.

On \textbf{calibration}: the reliability diagrams on the right plot each model's predicted confidence against its actual accuracy. A perfectly calibrated model would follow the diagonal. You can see that the transfer learning models are \textbf{over-confident} --- their confidence curves bow above the diagonal --- with Expected Calibration Error around 0.15 to 0.16. The custom CNN is paradoxically the best-calibrated individual model at ECE 0.11. However, the \textbf{three-class ensemble achieves ECE of approximately 0.05} --- averaging six diverse checkpoints smooths over-confidence considerably, making it the best-calibrated model overall.

Finally, the ensemble also allows us to \textbf{adjust the decision threshold} without retraining. We can set it at maximum recall for screening --- guaranteeing at least 99\% sensitivity --- or at the F1-optimal threshold for assisted diagnosis. Two clinical operating points from one model.
}

%% ============================================================
\section*{EXPLAINABILITY — Jakub (approx.\ 1 min 30 s)}
%% ============================================================

% ------ SLIDE 17 ------
\slidetag{17 — Grad-CAM: Explainability}
\timing{$\sim$30 s}

\jakub{Slide 17 — Grad-CAM: Explainability}{
Grad-CAM works by computing the gradient of the predicted class score with respect to the last convolutional layer, producing a heatmap that highlights which spatial regions most influenced the prediction. It requires no changes to the model architecture. For medical AI, this matters enormously: we need to verify that the model is attending to the lungs --- not to text labels, scan borders, or JPEG artefacts. It also helps build trust with clinicians. As you can see in the four examples from DenseNet121: true positives activate over opacity regions, true negatives show low diffuse activation over clear lung fields, and false negatives show subtle infiltrates that the model underweighted --- visually explaining why they were missed.
}

\vspace{0.4em}

% ------ SLIDE 18 ------
\slidetag{18 — Grad-CAM: ResNet18 vs.\ DenseNet121}
\timing{$\sim$30 s}

\jakub{Slide 18 — Grad-CAM: ResNet18 vs.\ DenseNet121}{
Comparing the two models side by side across true positives, true negatives, false positives, and false negatives: both attend to clinically plausible lung regions --- neither is latching onto irrelevant artefacts. However, DenseNet121 consistently produces more focused and spatially coherent activations. This aligns with its higher AUROC and is consistent with the dense connectivity architecture encouraging finer feature localisation.
}

\vspace{0.4em}

% ------ SLIDE 19 ------
\slidetag{19 — Grad-CAM: Three-Class Model}
\timing{$\sim$25 s}

\jakub{Slide 19 — Grad-CAM: Three-Class Model}{
Finally, here is Grad-CAM applied to the three-class model. The correct Bacteria prediction on the left shows focused activation over a focal consolidation. The correct Normal case in the centre shows low diffuse activation over clear lung fields. The most instructive case is the error on the right --- a Normal lung misclassified as Virus. The heatmap shows diffuse activation over an otherwise normal-looking lung, visually explaining the dominant error mode of the three-class task: viral infiltrates are so subtle and bilateral that the model cannot reliably distinguish them from normal texture.
}

%% ============================================================
\section*{ETHICS \& CONCLUSIONS — Carolina (approx.\ 1 min)}
%% ============================================================

% ------ SLIDE 20 ------
\slidetag{20 — Ethics \& Limitations}
\timing{$\sim$30 s}

\carol{Slide 20 — Ethics \& Limitations}{
We want to be transparent about the boundaries of this work. On the bias side: data comes from a single institution in Guangzhou, limiting generalisability. The dataset is pediatric only and may not transfer to adults or different equipment. Class imbalance creates a tendency to over-predict pneumonia, and JPEG artefacts or embedded text markers could be spuriously used as features. Technically: the validation set is only 16 images, there is no external validation cohort, and Grad-CAM is an interpretability tool, not ground truth. Most importantly --- this system is not suitable for clinical deployment. Real-world use would require multi-site validation, a calibration and bias audit, clinician oversight, and regulatory approval.
}

\vspace{0.4em}

% ------ SLIDE 21 ------
\slidetag{21 — Conclusions}
\timing{$\sim$30 s}

\carol{Slide 21 — Conclusions}{
To summarise: DenseNet121 is our best model, achieving AUROC of 0.966, recall of 0.997, and F1 of 0.887. Transfer learning substantially outperforms training from scratch across all metrics. All three success criteria were met. The ensemble misses just one pneumonia case out of 390, and inference cost is identical across all models at roughly 13 to 15 milliseconds. For future work, we see clear paths toward multi-institution and adult datasets, severity grading, better handling of class imbalance through weighted or focal loss, temperature scaling for calibration, and eventually prospective clinical evaluation. Thank you.
}

\vspace{0.4em}

% ------ SLIDES 22–24 (References) ------
\slidetag{22--24 — References}

\noindent\textit{[Skip --- references slides are displayed but not read aloud. If asked, you can point to the slide number.]}

\vspace{0.4em}

% ------ SLIDE 25 ------
\slidetag{25 — Thank You}

\noindent\textit{[Skip --- thank you slide is displayed but not read aloud. If asked, you can point to the slide number.]}


\noindent\rule{\linewidth}{1.2pt}

%% ============================================================
\newpage
\section*{PROJECT SUMMARY — For Complete Context}
%% ============================================================

\begin{tcolorbox}[sumbox={What This Project Is About — A Complete Overview}]
\textbf{The clinical problem.}
Pneumonia kills more than 2.5 million people per year and is the leading cause of death in children under five worldwide. Diagnosing it requires a trained radiologist to interpret a chest X-ray, but radiologists are scarce in low-income settings. Chest X-rays are the most accessible imaging modality globally, so automating their interpretation could save lives. This project asks: can a small, efficient deep learning model reliably detect pneumonia from chest X-rays, and can we understand what it is looking at?

\medskip
\textbf{The dataset.}
We used the Kermany et al.\ 2018 dataset: 5,856 pediatric chest X-rays from the Guangzhou Women and Children's Medical Centre, labelled as NORMAL or PNEUMONIA. The dataset is publicly available on Kaggle. A key property is that each pneumonia filename encodes whether the infection is bacterial or viral --- allowing us to derive three-class labels for free, without additional manual annotation. The training set is 74\% pneumonia (heavily imbalanced), which is why we optimised for recall and AUROC rather than raw accuracy.

\medskip
\textbf{Preprocessing.}
All images are resized to $224 \times 224$ pixels, converted from grayscale to RGB by triplicating the single channel, and normalised with ImageNet statistics. Training images additionally receive random horizontal flips and rotations up to $\pm 10$ degrees. Augmentation is deliberately conservative to avoid distorting clinically relevant features such as opacities.

\medskip
\textbf{The three models.}
We compared three architectures:
\begin{itemize}[nosep]
  \item \textbf{Custom CNN (baseline)}: three convolutional blocks, no pretrained weights, $\approx$0.5M parameters. Serves as a from-scratch lower bound.
  \item \textbf{ResNet18}: 18-layer residual network with skip connections, ImageNet-pretrained, $\approx$11M parameters. Transfer learning.
  \item \textbf{DenseNet121}: 121-layer densely-connected network where every layer receives feature maps from all preceding layers. ImageNet-pretrained, $\approx$7M parameters. Chosen because it was used by CheXNet (Rajpurkar et al.\ 2017) to reach radiologist-level pneumonia detection. Its dense feature reuse makes it especially good at detecting subtle diffuse opacities.
\end{itemize}

\medskip
\textbf{Training strategy.}
All models used the Adam optimiser ($lr=10^{-4}$), ReduceOnPlateau scheduling, batch size 32, early stopping with patience 5, and a maximum of 30 epochs. For the pretrained models, we applied a two-phase transfer learning strategy: Phase 1 trains only the classification head for 5 epochs with the backbone frozen; Phase 2 unfreezes the entire network for full fine-tuning. Every experiment was repeated with three random seeds, and we report mean $\pm$ standard deviation. The test set was evaluated exactly once at the end --- no data leakage.

\medskip
\textbf{Binary classification results.}
DenseNet121 is the best single model: recall $= 0.997 \pm 0.002$ (only 1--2 missed cases per run out of 390), AUROC $= 0.966 \pm 0.007$ (exceeding the CheXNet-14 reference of 0.927). All three models exceed the recall threshold of 0.90, but transfer learning models are substantially superior. Inference time is 13--15 ms per image for all three models --- no computational cost to choosing the better model.

\medskip
\textbf{Binary ensemble.}
Combining probability outputs from six model checkpoints (DenseNet121 and ResNet18, three seeds each) gives: F1 $= 0.896$, recall $= 0.997$, AUROC $= 0.963$, and \textbf{only 1 missed pneumonia case out of 390}. Specificity is 0.620, meaning 89 healthy patients are flagged --- acceptable in triage, since a clinician reviews flagged cases.

\medskip
\textbf{Innovation I: Three-class classification.}
We extended the task to distinguish Normal, Bacterial, and Viral pneumonia --- a clinically meaningful distinction since bacterial infections require antibiotics while viral ones do not. Using a patient-aware split to prevent leakage, the best single model (DenseNet121) achieves accuracy 0.799 and macro-F1 0.784. The ensemble of ResNet18 and DenseNet121 across three seeds achieves \textbf{accuracy 0.822, macro-F1 0.808, AUROC 0.960}. The hardest class is Virus (F1 0.72); the dominant error is Normal being misclassified as Virus because viral infiltrates can be subtle and diffuse.

\medskip
\textbf{Innovation II: Hierarchical two-stage classifier.}
Inspired by the Kermany et al.\ 2018 paper, which used two separate binary classifiers rather than a flat three-way model, we built a two-stage system: Stage A (DenseNet121) classifies Normal vs.\ Pneumonia; Stage B (DenseNet121, trained on pneumonia only) then classifies Bacteria vs.\ Virus. This hierarchical approach does \textbf{not} outperform the flat softmax (F1 macro 0.774 vs.\ 0.784). Reason: Stage A errors propagate to Stage B without recovery, and the Normal-to-Virus confusion is intrinsic to the data, not addressable by the architecture. This is a negative but informative result: the Virus bottleneck is a data-level problem.

\medskip
\textbf{Calibration.}
Transfer learning models are over-confident (Expected Calibration Error $\approx 0.15$--$0.16$); the custom CNN is paradoxically best-calibrated as a single model (ECE 0.11). The three-class ensemble achieves ECE $\approx 0.05$ --- averaging six diverse checkpoints smooths over-confidence, making it the best-calibrated model overall.

\medskip
\textbf{Grad-CAM explainability.}
Gradient-weighted Class Activation Mapping (Grad-CAM) produces a heatmap over the X-ray by computing gradients of the predicted class score with respect to the last convolutional layer's activations. For true positives, the heatmap activates over opacity regions. For true negatives, activation is low and diffuse over clear lung fields. For false negatives, the missed subtle infiltrate is weakly activated. DenseNet121 consistently produces more focused, spatially coherent heatmaps than ResNet18, consistent with its higher AUROC. For the three-class model, the Normal$\to$Virus error mode is visually explained: the model activates diffusely over an otherwise normal lung, treating faint bilateral texture as viral infiltrate.

\medskip
\textbf{Ethics and limitations.}
The dataset comes from a single institution (Guangzhou), limiting generalisability. It is pediatric only. Class imbalance biases predictions toward pneumonia. Validation split is only 16 images. No external validation cohort. Grad-CAM is not ground-truth localisation. \textbf{This system is not suitable for clinical deployment} and is presented as a research prototype. Real deployment would require multi-site validation, calibration and bias audit, clinician oversight, and regulatory approval.

\medskip
\textbf{Key numbers to remember:}
\begin{itemize}[nosep]
  \item Binary ensemble: recall $= 0.997$, F1 $= 0.896$, AUROC $= 0.963$, 1 missed case / 390
  \item Three-class ensemble: macro-F1 $= 0.808$, AUROC $= 0.960$
  \item DenseNet121 AUROC $= 0.966$ (CheXNet-14 reference: 0.927)
  \item Inference: $\sim$13--15 ms / image for all three models
  \item Calibration ECE: single models 0.11--0.16; three-class ensemble 0.05
\end{itemize}
\end{tcolorbox}

%% ============================================================
\newpage
\section*{GLOSSARY}
%% ============================================================

\begin{tcolorbox}[glossbox={Key Terms Explained}]

\textbf{Accuracy} --- The fraction of all predictions that are correct. Misleading when classes are imbalanced; a model that predicts ``pneumonia'' for every image would achieve 74\% accuracy on this dataset.

\medskip
\textbf{AUROC (Area Under the Receiver Operating Characteristic Curve)} --- A threshold-independent metric ranging from 0 to 1. It measures the probability that the model ranks a randomly chosen positive case above a randomly chosen negative case. 1.0 is perfect; 0.5 is random chance. More informative than accuracy for imbalanced datasets.

\medskip
\textbf{Batch Normalisation (BN)} --- A technique applied after each convolutional layer that normalises activations within a mini-batch. Stabilises and accelerates training.

\medskip
\textbf{Calibration / ECE (Expected Calibration Error)} --- A calibrated model is one whose confidence scores match its actual accuracy. ECE measures the average gap between confidence and accuracy across probability bins. ECE $= 0$ is perfect calibration; lower is better.

\medskip
\textbf{Chest X-ray} --- A 2D projection image created by passing X-rays through the body. The lungs appear as dark regions (air is radiolucent); opacities (white areas within the lungs) indicate fluid, infection, or consolidation.

\medskip
\textbf{CNN (Convolutional Neural Network)} --- A neural network architecture designed for image processing. Convolutional layers learn local spatial filters (e.g., edges, textures); pooling layers reduce spatial resolution; fully connected layers produce the final prediction.

\medskip
\textbf{Consolidation} --- A radiological finding where the air in the lung's alveoli is replaced by fluid, pus, or inflammatory cells. Appears as an opacity on X-ray. Classic sign of bacterial pneumonia.

\medskip
\textbf{DenseNet121 (Densely Connected Convolutional Network)} --- A 121-layer deep network where each layer receives as input the feature maps from \emph{all} preceding layers (not just the previous one). This dense connectivity promotes feature reuse and gradient flow, allowing a smaller number of parameters (7M) to learn richer representations.

\medskip
\textbf{Ensemble} --- Combining predictions from multiple independently trained models to reduce variance and improve robustness. Here, we average the probability outputs of six checkpoints (DenseNet121 and ResNet18 across three seeds).

\medskip
\textbf{F1 Score} --- The harmonic mean of precision and recall: $F1 = 2 \cdot \frac{\text{Precision} \cdot \text{Recall}}{\text{Precision} + \text{Recall}}$. Balances both metrics; useful when both false positives and false negatives matter.

\medskip
\textbf{False Negative (FN)} --- A pneumonia case that the model classified as Normal (a missed diagnosis). In a clinical context, this is the most dangerous error type.

\medskip
\textbf{False Positive (FP)} --- A healthy patient flagged as having pneumonia. In triage, this is acceptable because the case is reviewed by a clinician. Fewer FP means higher specificity.

\medskip
\textbf{Fine-tuning} --- After loading ImageNet-pretrained weights, we continue training the model (Phase 2) on the chest X-ray task with a low learning rate. This adapts pretrained features to the target domain.

\medskip
\textbf{Grad-CAM (Gradient-weighted Class Activation Mapping)} --- A visualisation method that produces a heatmap over the input image showing which regions most influenced the model's prediction. Computed by taking the gradient of the predicted class score with respect to the last convolutional layer's activations, then weighting and globally pooling those activations.

\medskip
\textbf{ImageNet} --- A large image database with over 14 million images and 1,000 classes (dogs, cats, vehicles, etc.). Models pretrained on ImageNet have learned general visual features (edges, textures, shapes) that transfer well to other domains including medical imaging.

\medskip
\textbf{Infiltrate} --- A non-specific radiological term for an opacity in the lung that can indicate infection, inflammation, or fluid. Viral pneumonia often presents as bilateral diffuse infiltrates.

\medskip
\textbf{Learning Rate Scheduler (ReduceOnPlateau)} --- Automatically reduces the learning rate by a factor (here 0.5) when the validation metric stops improving for a defined number of epochs (here patience = 3). Allows fine-grained convergence without manual tuning.

\medskip
\textbf{Macro-F1} --- The average of per-class F1 scores, treating all classes equally regardless of size. Appropriate for evaluating multi-class performance when class sizes differ.

\medskip
\textbf{Opacity} --- Any region on an X-ray that appears whiter than expected. In the lungs, opacities indicate that air has been partially or fully replaced by denser material (fluid, pus, tumour).

\medskip
\textbf{Patient-aware split} --- A data splitting strategy ensuring that all X-rays from the same patient belong to the same split (train, validation, or test). Prevents a form of data leakage where the model could partially memorise patient-specific anatomy rather than learning the classification task.

\medskip
\textbf{Precision} --- Of all cases the model predicted as pneumonia, what fraction actually had pneumonia? High precision means few false alarms. $\text{Precision} = \frac{TP}{TP + FP}$.

\medskip
\textbf{Recall (Sensitivity)} --- Of all actual pneumonia cases, what fraction did the model correctly detect? High recall means few missed cases. $\text{Recall} = \frac{TP}{TP + FN}$. Our primary success criterion (target: $> 0.90$).

\medskip
\textbf{ResNet18 (Residual Network, 18 layers)} --- A convolutional network that introduces \emph{skip connections}: a shortcut that adds the input of a block directly to its output. Prevents the vanishing gradient problem and allows training of much deeper networks.

\medskip
\textbf{Sigmoid} --- An activation function that maps any real number to $(0, 1)$, used as the output for binary classification to produce a probability.

\medskip
\textbf{Softmax} --- A generalisation of sigmoid for multi-class classification. Converts a vector of raw scores into a probability distribution over classes that sums to 1.

\medskip
\textbf{Specificity} --- Of all healthy (Normal) patients, what fraction did the model correctly identify as Normal? $\text{Specificity} = \frac{TN}{TN + FP}$. Low specificity means many false alarms.

\medskip
\textbf{Transfer Learning} --- Using a model pretrained on one large dataset (ImageNet) as a starting point for a different task (chest X-ray classification). The pretrained features provide a strong initialisation, dramatically reducing the amount of labelled medical data needed.

\medskip
\textbf{True Negative (TN)} --- A healthy patient correctly classified as Normal.

\medskip
\textbf{True Positive (TP)} --- A pneumonia patient correctly classified as Pneumonia.

\medskip
\textbf{Two-phase transfer learning} --- Phase 1: freeze the pretrained backbone; train only the new classification head for a few epochs to stabilise it. Phase 2: unfreeze the entire network and fine-tune jointly at a low learning rate.

\end{tcolorbox}

%% ============================================================
\newpage
\section*{PROBABLE Q\&A}
%% ============================================================

\begin{tcolorbox}[qabox={Q1 — Why is specificity so low (0.62 for the ensemble)?}]
\textbf{Answer:} Specificity and recall are in direct tension. When we maximise recall to near 1.0, we are pushing the decision boundary to capture almost every positive case --- and that inevitably means accepting more false positives from the healthy class. Our clinical framing is triage: the system flags suspicious cases for clinician review rather than making autonomous decisions. In triage, 89 false alarms out of 234 healthy patients is acceptable because each flag gets reviewed by a human. What is not acceptable is missing true pneumonia cases, which we virtually eliminate (1 miss out of 390). If specificity mattered more, we could simply lower the decision threshold --- no retraining needed.
\end{tcolorbox}

\begin{tcolorbox}[qabox={Q2 — Why did you choose DenseNet121 as your primary model rather than a more recent architecture like EfficientNet or Vision Transformers?}]
\textbf{Answer:} The project brief was ``lightweight deep learning models'', and we wanted an architecture with a direct clinical precedent. DenseNet121's CheXNet heritage --- Rajpurkar et al.\ 2017 demonstrating radiologist-level performance --- made it the natural choice for a chest X-ray task. It also achieves more parameters than the custom CNN but fewer than ResNet18 (7M vs.\ 11M), while outperforming ResNet18 on AUROC. We did not explore ViTs or EfficientNet due to scope constraints, but those would be natural extensions. Interestingly, DenseNet121 already exceeded the CheXNet-14 AUROC reference (0.966 vs.\ 0.927), suggesting the pretrained weights and fine-tuning strategy were effective.
\end{tcolorbox}

\begin{tcolorbox}[qabox={Q3 — Why is Virus so much harder to classify than Bacteria?}]
\textbf{Answer:} Viral pneumonia tends to produce \emph{bilateral diffuse infiltrates} --- subtle, widespread, low-density changes across both lung fields. These can look very similar to a normal lung on a standard anterior-posterior X-ray, especially in early stages or in paediatric patients where lung volumes are smaller. Bacterial pneumonia, in contrast, typically produces \emph{focal lobar consolidation} --- a denser, well-defined opacity usually confined to one lobe --- which is visually much more distinctive. This is confirmed by our Grad-CAM visualisations: the false Normal-to-Virus cases show diffuse, non-specific activation over what visually appears to be a clear lung. The bottleneck is in the data itself, not the architecture.
\end{tcolorbox}

\begin{tcolorbox}[qabox={Q4 — Why did the hierarchical classifier not improve over the flat softmax?}]
\textbf{Answer:} Two reasons. First, \textbf{error propagation}: any mistake in Stage A (Normal vs.\ Pneumonia) cannot be corrected in Stage B. In the flat model, a softmax over three classes can implicitly hedge between Normal and Virus even when uncertain. In the cascade, if Stage A commits to Pneumonia, Stage B never has a chance to recover and say ``actually this looks more like Normal.'' Second, the \textbf{dominant error mode} --- Normal misclassified as Virus --- exists at the feature level, not the architectural level. Both stages of the hierarchy see the same subtle visual features and make the same mistake. The flat softmax weighs all three classes jointly and is empirically more robust. This result is important because it replicates and questions the original paper's design choice.
\end{tcolorbox}

\begin{tcolorbox}[qabox={Q5 — The validation set has only 16 images. Is that a concern?}]
\textbf{Answer:} Yes, and we acknowledge it explicitly in our limitations slide. A 16-image validation set provides very noisy estimates of generalisation performance during training. In practice, we used it only for early stopping and learning rate scheduling --- decisions that are relatively robust to noise --- and all final hyperparameters were fixed before any test-set evaluation. The test set (624 images) is much larger and was evaluated only once. We also mitigated this by running three random seeds and reporting mean $\pm$ std, which gives a measure of robustness. Ideally, a larger validation set or $k$-fold cross-validation would be preferred, but those were constraints of the Kaggle dataset split we used.
\end{tcolorbox}

\begin{tcolorbox}[qabox={Q6 — How does Grad-CAM actually work technically?}]
\textbf{Answer:} Grad-CAM works by: (1) performing a forward pass and obtaining the predicted class score; (2) computing the gradient of that score with respect to the \emph{feature maps} of the last convolutional layer --- this tells us how much each spatial location in that feature map contributed to the final prediction; (3) globally average-pooling those gradients over the spatial dimensions to get a weight per feature map channel; (4) computing a weighted sum of the feature maps using those weights; (5) applying a ReLU to keep only positive contributions; and (6) upsampling the resulting heat map to the original image size and overlaying it. The result is a coarse localisation map that highlights which image regions were most discriminative for the predicted class. It requires no modification to the network and works with any architecture that has convolutional layers.
\end{tcolorbox}

\begin{tcolorbox}[qabox={Q7 — Could this system work in a real clinical setting?}]
\textbf{Answer:} Not in its current form, and we are explicit about this. The limitations are substantial: single-institution training data from one paediatric centre in China limits geographical and demographic generalisability; there is no external validation cohort; the validation split is too small; and we have not conducted the regulatory, clinical, or ethical review required for medical device certification. What this system demonstrates is that \emph{the classification signal exists} and is learnable with lightweight, open-source models. The path to clinical deployment would require multi-site prospective studies, independent validation on adult data and different imaging equipment, uncertainty quantification, clinician-in-the-loop integration, and regulatory approval under frameworks such as the EU Medical Device Regulation or FDA clearance. The model weights are publicly available on Hugging Face for further research.
\end{tcolorbox}

\begin{tcolorbox}[qabox={Q8 — Why do the transfer learning models have higher ECE (worse calibration) than the custom CNN?}]
\textbf{Answer:} This is a known phenomenon in transfer learning. Pretrained models often produce high-confidence outputs because they have been trained on very large datasets where confident predictions are rewarded. When fine-tuned on a small target dataset, the confidence can remain artificially high even when the model is less certain. The custom CNN, trained from scratch on a limited dataset, naturally produces more modest confidence scores --- it has not been ``taught to be confident'' by large-scale pretraining. However, this does not mean the CNN is a better model --- its discriminative performance (AUROC, recall) is substantially worse. The ensemble mitigates over-confidence by averaging six diverse checkpoints, bringing ECE down to approximately 0.05 for the three-class case. Temperature scaling would be the next step to further improve calibration without changing the model architecture.
\end{tcolorbox}

\end{document}
%% ============================================================
